\subsection{紧线性映射与弱拓扑}

设 E 是局部凸空间，\(E'\) 是它的对偶。回顾，E 上的拓扑 \(\sigma(\mathrm{E},\mathrm{E}')\) 称为 E 上的弱拓扑（TVS, IV, p. 4）。在本节中，\(E_{\sigma}\) 表示赋有弱拓扑的空间 E。

命题 6. --- 设 E 是局部凸空间，F 是分离的局部凸空间，\(u: \mathrm{E} \to \mathrm{F}\) 是紧线性映射，B 是 E 的一个有界子集。\(u\) 在 B 上的限制是从赋有 \(\sigma(\mathrm{E}, \mathrm{E}')\) 诱导拓扑的集合 B 到空间 F 的连续映射。

映射 \(u\) 是从 E 到 F 的连续线性映射，也是从 \(\mathrm{E}_{\sigma}\) 到 \(\mathrm{F}_{\sigma}\) 的连续线性映射。因此，它在 B 上的限制是从赋有 \(\sigma(\mathrm{E}, \mathrm{E}')\) 诱导拓扑的集合 B 到空间 \(\mathrm{F}_{\sigma}\) 的连续映射。而集合 \(u(\mathrm{B})\) 包含于 F 的一个紧子集 C（III, p. 2, 注记 1），且空间 \(\mathrm{F}_{\sigma}\) 是豪斯多夫的，故由 F 与由 \(\mathrm{F}_{\sigma}\) 在 C 上诱导的拓扑重合。命题得证。

推论. --- 设 \((x_{n})_{n\in\mathbf{N}}\) 是 E 的点的一个序列，它对弱拓扑收敛于 E 的一个点 x。则序列 \((u(x_{n}))_{n\in\mathbf{N}}\) 在 F 中收敛于 \(u(x)\)。

事实上，由点 \(x\) 与序列 \((x_{n})_{n\in \mathbf{N}}\) 的诸点所组成的集是 E 的一个有界子集（EVT, III, p. 3, 命题 2 之推论）。

命题 7. --- 设 E 是半赋范空间，E' 是它的对偶，B 是 E' 的单位球，u 是从 E' 到一个分离局部凸空间 F 的线性映射。若 u 在赋有 \(\sigma(\mathrm{E}',\mathrm{E})\) 诱导拓扑的 B 上的限制是连续的，则 \(u(\mathrm{B})\) 是 F 的一个紧子集，且 u 是从 E 的强对偶 \(\mathrm{E}_b'\) 到 F 的紧线性映射。

事实上，赋有 \(\sigma(\mathrm{E}^{\prime},\mathrm{E})\) 诱导拓扑的集合 B 是紧的（EVT, III, p. 17, cor. 2）。

推论 1. --- 设 E 是可数型的半赋范空间，E' 是它的对偶，B 是 E' 的单位球。设 u 是从 E' 到一个分离局部凸空间 F 的线性映射。假设对 B 中元素的每个序列 \((x_{n})_{n\in\mathbf{N}}\)（它对 \(\sigma(\mathrm{E}',\mathrm{E})\) 收敛于 0），序列 \((u(x_{n}))_{n\in\mathbf{N}}\) 都在 F 中收敛于 0。那么，\(u(\mathrm{B})\) 是 F 的一个紧子集，且 u 是从 E 的强对偶 \(E_{b}'\) 到 F 的紧线性映射。

事实上，\(\sigma(\mathrm{E}^{\prime},\mathrm{E})\) 在 B 上诱导的拓扑是可度量化的（EVT, III, p. 19, cor. 2）。因此，推论的假设蕴含：\(u\) 在 B 上的限制（当 B 赋有拓扑 \(\sigma(\mathrm{E}^{\prime},\mathrm{E})\) 时）是从 B 到 F 的连续映射，于是推论由前一命题得出。

推论 2. --- 设 E 是半赋范空间，S 是 E 的预紧子集所成的一个集，其并在 E 中稠密。用 \(E'_{b}\) 表示 E 的强对偶，用 \(E'_{S}\) 表示赋有 S-拓扑的 E 的对偶。空间 \(E'_{S}\) 是分离的，且从 \(E'_{b}\) 到 \(E'_{S}\) 的恒等映射是紧的。

依照 TG, X, p. 8, prop. 7，空间 \(E_{S}^{\prime}\) 是分离的；在 \(E^{\prime}\) 的单位球上，S-拓扑与弱拓扑 \(\sigma(\mathrm{E}^{\prime},\mathrm{E})\) 重合（EVT, III, p. 17, prop. 5），由此即得推论。

例. --- 设 E 是半赋范空间。从 \(E_{b}^{\prime}\) 到赋有弱拓扑、紧收敛拓扑或预紧收敛拓扑的空间 \(E^{\prime}\) 的恒等映射是紧的。

命题 8. --- 设 E 是自反的巴拿赫空间，B 是它的单位球，F 是分离的局部凸空间，u: E → F 是线性映射。下列条件等价：

\begin{enumerate}
\def\labelenumi{(\roman{enumi})}
\item
  线性映射 u 是紧的；
\item
  集合 \(u(\mathrm{B})\) 是 F 的一个紧子集；
\item
  u 在 B 上的限制是从赋有 \(\sigma(\mathrm{E},\mathrm{E}')\) 诱导拓扑的集合 B 到 F 的连续映射；
\item
  对 B 中点的任何序列 \((x_{n})_{n\in \mathbf{N}}\)（它对拓扑 \(\sigma (\mathrm{E},\mathrm{E}^{\prime})\) 收敛于 0），序列 \((u(x_n))_{n\in \mathbf{N}}\) 在 F 中收敛于 0；
\item
  从 \(u(\mathrm{B})\) 的点的任何无穷序列中都能取出一个在 F 中收敛的子序列；
\item
  \(u(\mathrm{B})\) 的点的任何无穷序列在 F 中都有聚点。
\end{enumerate}

由于 E 是自反空间，B 是 \(E_{\sigma}\) 的一个紧子集（EVT, IV, p. 17, prop. 6），故 (iii) 蕴含 (ii)。按定义，条件 (ii) 蕴含 (i)，而由命题 6，条件 (i) 蕴含 (iii)。同样初等的是 (iii) 蕴含 (iv) 且 (v) 蕴含 (vi)。

我们现在证明 (iv) 蕴含 (v)。设 \((x_{n})\) 是 B 的点的一个无穷序列。由于 B 是 \(E_{\sigma}\) 的紧子集，由什穆利安定理（EVT, IV, p. 36, th. 2），存在一个序列 \((y_{n})\)（从序列 \((x_{n})\) 中取出），它在 \(E_{\sigma}\) 中收敛于一个极限 y。序列 \((y_{n}-y)\) 在 \(E_{\sigma}\) 中有界，从而在 E 中有界（EVT, IV, p. 1, prop. 1）；因此它包含于 B 的一个位似集合。若条件 (iv) 成立，则序列 \((u(y_{n}-y))\) 在 F 中收敛于 0，而 \((u(y_{n}))\) 作为 \((u(x_{n}))\) 的一个子序列收敛于 \(u(y)\)。这就证明了 (iv) 蕴含 (v)。

最后，我们证明 (vi) 蕴含 (ii)。设 \(j: \mathrm{F} \to \widehat{\mathrm{F}}\) 是 F 到其完备化中的典范单射。在假设 (vi) 之下，\(u(\mathrm{B})\) 是 F 的一个预紧子集（EVT, IV, p. 32, prop. 1），因此 \(j(u(\mathrm{B}))\) 是 \(\widehat{\mathrm{F}}\) 的一个相对紧子集（TG, II, 第 2 目, p. 29），且 \(j \circ u\) 是紧线性映射。由已经证明的条件 (i) 与 (ii) 的等价性，\(j(u(\mathrm{B}))\) 是 \(\widehat{\mathrm{F}}\) 的一个紧子集。因此，\(u(\mathrm{B})\) 是 F 的一个紧子集。

